\section{Integral and Gaussian forms of the four-distance equations}
\label{sec:integral}

We label the vertices of the unit square by
\[
 P_0=(0,0),\qquad P_1=(1,0),\qquad
 P_2=(0,1),\qquad P_3=(1,1).
\]
A \emph{strict rational-distance point} is a point $(x,y)$ with
$0<x,y<1$ whose distances $r_i$ to all four $P_i$ belong to
$\Q_{>0}$. Thus
\begin{equation}\label{eq:int:four}
\begin{aligned}
 r_0^2&=x^2+y^2,&r_1^2&=(1-x)^2+y^2,\\
 r_2^2&=x^2+(1-y)^2,&r_3^2&=(1-x)^2+(1-y)^2.
\end{aligned}
\end{equation}
The labels in these equations will be retained throughout.

\subsection{Recovering the point from its four radii}

Subtracting adjacent equations gives
\begin{equation}\label{eq:int:recover}
 2x=1+r_0^2-r_1^2,\qquad
 2y=1+r_0^2-r_2^2.
\end{equation}
In particular the coordinates are rational. The following elimination
keeps both the point and its four positive radii.

\begin{proposition}\label{prop:int:radius-model}
Positive rational numbers $(r_0,r_1,r_2,r_3)$ arise from a unique strict
rational-distance point if and only if
\begin{align}
 r_0^2+r_3^2&=r_1^2+r_2^2,\label{eq:int:flag}\\
 (1+r_0^2-r_1^2)^2+(1+r_0^2-r_2^2)^2&=4r_0^2,\label{eq:int:radiusquartic}\\
 -1<r_0^2-r_1^2&<1,\qquad
 -1<r_0^2-r_2^2<1.\label{eq:int:strict-radii}
\end{align}
The point is given by \eqref{eq:int:recover}.
\end{proposition}

\begin{proof}
Equations \eqref{eq:int:flag} and \eqref{eq:int:radiusquartic} follow
by adding and subtracting \eqref{eq:int:four}. Conversely, define $x,y$
by \eqref{eq:int:recover}. The inequalities give $0<x,y<1$, and
\eqref{eq:int:radiusquartic} gives $x^2+y^2=r_0^2$. Subtracting $2x-1$
and $2y-1$ gives the equations for $r_1,r_2$; the remaining equation
then follows from \eqref{eq:int:flag}. Positivity fixes the four square
roots, and \eqref{eq:int:recover} gives uniqueness.
\end{proof}

Let $L$ be the least positive common denominator of $x,y$, and put
\begin{equation}\label{eq:int:lattice}
 X=Lx,\qquad Y=Ly,\qquad A_i=Lr_i.
\end{equation}
Then $X,Y,L$ are integers with
\[
 0<X,Y<L,\qquad \gcd(X,Y,L)=1.
\]
Each $A_i$ is an integer: a rational number whose square is an integer
is itself an integer. Equations \eqref{eq:int:four} become
\begin{equation}\label{eq:int:integer-four}
\begin{aligned}
 A_0^2&=X^2+Y^2,&A_1^2&=(L-X)^2+Y^2,\\
 A_2^2&=X^2+(L-Y)^2,&A_3^2&=(L-X)^2+(L-Y)^2.
\end{aligned}
\end{equation}
Conversely, positive integers satisfying these equations and
$0<X,Y<L$ give a strict rational-distance point by division by $L$;
the additional gcd condition makes $L$ its least coordinate denominator.

\subsection{Parity, symmetry lines, and the primitive denominator}

We first record the elementary descent needed to deal with the symmetry
lines of the square.

\begin{lemma}[Square-area descent]\label{lem:int:fermat}
A rational right triangle cannot have positive rational-square area.
\end{lemma}

\begin{proof}
Clearing denominators would give an integer right triangle of square
area. Dividing by the gcd of its sides preserves this property: the
new area is an integer and a rational square, hence an integer square.
Choose a primitive such triangle with least hypotenuse. Its legs and
hypotenuse have the form
\[
 2rs,\quad r^2-s^2,\quad r^2+s^2,
 \qquad r>s>0,\quad \gcd(r,s)=1,
\]
where $r,s$ have opposite parity. To obtain this parametrization, write
the odd leg as $a$ and the hypotenuse as $h$. The coprime positive
integers $(h+a)/2,(h-a)/2$ have square product and are therefore
squares.

The area $rs(r-s)(r+s)$ is a square. Its four factors are pairwise
coprime, so
\[
 r=u^2,\qquad s=v^2,\qquad r+s=w^2,\qquad r-s=z^2.
\]
The numbers $w,z$ are odd. Since $w^2-z^2=2v^2$ and odd squares are
$1$ modulo $8$, the integer $v$ is even. The positive integers
\[
 a'=(w+z)/2,\qquad b'=(w-z)/2
\]
satisfy
\[
 (a')^2+(b')^2=u^2,\qquad a'b'/2=(v/2)^2.
\]
They give a right triangle of square area with hypotenuse $u$ strictly
smaller than $r^2+s^2$. Primitive reduction, if required, decreases
the hypotenuse further. This contradicts its minimal choice.
\end{proof}

\begin{lemma}\label{lem:int:symmetry}
A strict rational-distance point lies on neither a midline nor a
diagonal of the square. Its four radii are pairwise distinct.
\end{lemma}

\begin{proof}
On $y=1/2$, put $u=2x-1$, $p=2r_0$, and $q=2r_1$. Then
\[
 p^2=u^2+2u+2,\qquad q^2=u^2-2u+2,
 \qquad (pq)^2=u^4+4.
\]
If $u\ne0$, the rational numbers $|u|,2/|u|,pq/|u|$ are the sides
of a right triangle of area one, contrary to
Lemma~\ref{lem:int:fermat}. If $u=0$, the squared distance to a vertex
is $1/2$, which is not a rational square. The other midline follows
by exchanging coordinates. On $x=y$ the distance to $P_0$ is
$\sqrt2\,x$; on $x+y=1$ the distance to $P_1$ is
$\sqrt2\,(1-x)$. These are irrational because the coordinates are
rational and the indicated factors are positive. Finally, the
perpendicular bisector of every pair of square vertices is one of
these four lines.
\end{proof}

\begin{proposition}[Primitive congruences]\label{int:primitive}
For every strict rational-distance point,
\[
 24\mid L.
\]
After a single exchange of the coordinates, $X$ is odd and $4\mid Y$.
All four $A_i$ are odd and prime to $3$. Every prime divisor of
\begin{equation}\label{eq:int:diagonal-product}
 (X-Y)(X+Y-L)
\end{equation}
is congruent to $1$ or $7$ modulo $8$.
\end{proposition}

\begin{proof}
If $L$ were odd, one could choose an odd element from each of
$\{X,X-L\}$ and $\{Y,Y-L\}$. The corresponding equation in
\eqref{eq:int:integer-four} would have right side $2$ modulo $4$.
Thus $L$ is even. Primitivity excludes both $X,Y$ even, and their
squared sum excludes both odd. Exchange coordinates so that $X$ is
odd. All $A_i$ are odd; reduction modulo $8$ then gives $4\mid Y$
and $4\mid L-Y$, hence $4\mid L$.

If $3\nmid L$, each of $\{X,X-L\}$ and $\{Y,Y-L\}$ contains a
unit modulo $3$. The associated squared distance would be $2$ modulo
$3$. Hence $3\mid L$. Now exactly one of $X,Y$ is divisible by $3$:
both units would again give a nonsquare sum, while both nonunits would
contradict primitivity. Consequently every $A_i$ is a $3$-adic unit.

The integer in \eqref{eq:int:diagonal-product} is nonzero by
Lemma~\ref{lem:int:symmetry}, and both factors are odd. Let an odd
prime $p$ divide $X-Y$. If $X,Y$ are units modulo $p$, then $A_0^2$
reduces to twice a nonzero square. Otherwise both are zero modulo $p$;
primitivity gives $p\nmid L$, and $A_3^2$ reduces to $2L^2$.
In either case $\leg2p=1$. If instead $p\mid X+Y-L$, then
$A_1^2\equiv2Y^2$ and $A_2^2\equiv2X^2\pmod p$; at least one of
$X,Y$ is a unit, again by primitivity. Thus $\leg2p=1$ in this case
as well. The supplementary law for $2$ gives $p\equiv1,7\pmod8$.

A nonzero integer supported on these primes, with either sign, is
$1$ or $7$ modulo $8$. Were $L\equiv4\pmod8$, our orientation would
instead give
\[
 (X-Y)(X+Y-L)=X^2-Y^2-LX+LY\equiv5\pmod8.
\]
This contradiction proves $8\mid L$, and therefore $24\mid L$.
\end{proof}

\subsection{Common content of the integer radii}

The common divisor of the four radii has a precise relation to the
coordinate denominator.

\begin{proposition}\label{int:content}
Set
\[
 G=\gcd(A_0,A_1,A_2,A_3),\qquad a_i=A_i/G.
\]
Then $G^2\mid L$. Every prime divisor of $G$ is $1$ modulo $4$,
and $\gcd(G,6)=1$. Writing
\begin{equation}\label{eq:int:content-normalization}
 H=L/G^2,
\end{equation}
one has
\[
 24\mid H,\qquad r_i=\frac{a_i}{GH},\qquad
 \gcd(a_0,a_1,a_2,a_3)=1.
\]
For every prime $p\mid H$, all four $a_i$ are $p$-adic units.
\end{proposition}

\begin{proof}
The congruences already proved give $\gcd(G,6)=1$. Suppose an odd
prime $p$ divides $G$ but not $L$. Subtracting the origin square from
the two adjacent squares gives $X\equiv Y\equiv L/2\pmod p$.
Then $A_0^2\equiv L^2/2\ne0\pmod p$, a contradiction. Thus
$p\mid L$. Primitivity and $X^2+Y^2\equiv0\pmod p$ now force both
$X,Y$ to be units and $(Y/X)^2\equiv-1\pmod p$. In particular
$p\equiv1\pmod4$.

Write $s=v_p(G)$ and $e=v_p(L)$. The identity
\[
 A_1^2-A_0^2=L(L-2X)
\]
has right side of exact valuation $e$, since $p\nmid X$. Its left
side is divisible by $p^{2s}$; hence $e\ge2s$. This proves
$G^2\mid L$ at every prime. Division of the adjacent squared
differences by $G^2$ shows that all four $a_i^2$ have the same residue
modulo each prime divisor of $H$. If one were divisible by that prime,
all would be, contrary to their collective gcd. The remaining
assertions follow from \eqref{eq:int:content-normalization} and
$24\mid L$.
\end{proof}

There is also a uniform relation between the denominators of a sum
and a difference of two radii. Here the denominator of a nonzero
rational number always means its positive reduced denominator.

\begin{theorem}[Complementary denominators]\label{thm:int:pair-denominators}
For any distinct $i,j$, let $d_+$ and $d_-$ be the denominators of
$r_i+r_j$ and $r_i-r_j$. Then
\begin{equation}\label{eq:int:denominator-product}
 \gcd(d_+,d_-)=G,\qquad d_+d_-=L/2.
\end{equation}
\end{theorem}

\begin{proof}
Both denominators divide $L$. Fix an odd $p\mid L$, and write
$e=v_p(L)$, $s=v_p(G)$. If $s=0$, the four $A_i$ are units modulo
$p$, because their squares have the same residue and at least one is
a unit. Every difference $A_i^2-A_j^2$ is divisible by $p^e$.
Exactly one of $A_i+A_j,A_i-A_j$ is divisible by $p$, so their
reduced denominator depths are $0$ and $e$.

Suppose $s>0$. The proof of Proposition~\ref{int:content} shows
$Y/X\equiv\pm\sqrt{-1}\pmod p$. Thus
$L-2X$, $L-2Y$, $X-Y$, and $X+Y-L$ are all units modulo $p$.
The six squared differences are, up to sign,
\begin{equation}\label{eq:int:six-differences}
 L(L-2X),\quad L(L-2Y),\quad
 2L(X+Y-L),\quad 2L(X-Y),
\end{equation}
with repetitions in the parallel directions. Each has exact depth
$e$. If $e>2s$, the four $a_i$ are units, and their sum and difference
have depths $0$ and $e-2s$. The depths of the corresponding
denominators are consequently $e-s$ and $s$. If $e=2s$, the squared
difference of the normalized pair is a unit, so both its sum and
difference are units; both denominator depths are $s$. In all cases
the sum of the two depths is $e$ and their minimum is $s$.

At $2$, write $e=v_2(L)\ge3$ and keep $X$ odd, $4\mid Y$. All four
$A_i$ are odd. The horizontal and opposite squared differences in
\eqref{eq:int:six-differences} have depth $e+1$; the vertical ones
have depth at least $e+3$. Of a sum and a difference of two odd
integers, precisely one has depth one. The other therefore has depth
at least $e$. Their reduced denominator depths are $e-1$ and $0$.
Combining the prime-by-prime statements proves
\eqref{eq:int:denominator-product}.
\end{proof}

\subsection{Gaussian roots and the four corner classes}

Put $c_0=0$, $c_1=1$, $c_2=i$, $c_3=1+i$, and define the integer
Gaussian displacements
\[
 U_j=X+iY-Lc_j.
\]
For $j=(\epsilon,\eta)$ put
\[
 d_j=\gcd(|X-\epsilon L|,|Y-\eta L|),\qquad
 m_j=\sf(d_j),
\]
where $\sf$ denotes the positive squarefree part of a positive integer.
The coordinate orientation is the one in Proposition~\ref{int:primitive}.

\begin{theorem}[Canonical corner classes]\label{int:gaussian}
The integers $d_j$ are odd, pairwise coprime, and coprime to $L$.
Every prime dividing any $d_j$ is $1$ or $7$ modulo $8$. There are
nonzero $q_j\in\Z[i]$ such that
\begin{equation}\label{eq:int:gaussian-roots}
 U_j=m_jq_j^2,\qquad \Nm(q_j)=A_j/m_j>0.
\end{equation}
In particular $m_0,m_1,m_2,m_3$ are positive odd squarefree integers
with disjoint supports, and, for $z=x+iy$,
\begin{equation}\label{eq:int:gaussian-classes}
 [z-c_j]=[Lm_j]\quad\text{in }\Q(i)^*/\Q(i)^{*2}.
\end{equation}
Conversely, if a rational strict point has a primitive representation
with $X$ odd and admits \eqref{eq:int:gaussian-roots} with positive
integers $m_j$ and $q_j\in\Z[i]$, all four of its distances are rational.
\end{theorem}

\begin{proof}
Every first coordinate $X-\epsilon L$ is odd, so all $d_j$ are odd.
A prime dividing $d_j$ and $L$ would divide $X,Y,L$. A prime dividing
two different $d_j$ would, by subtracting their coordinate pairs,
divide $L$ as well. This proves the asserted coprimalities.
If $p\mid d_j$, the squared distance from the opposite vertex is
$2L^2$ modulo $p$. Since $p\nmid L$, this implies $\leg2p=1$.

Divide the two legs and hypotenuse of the $j$th Pythagorean triangle
by $d_j$. This gives coprime integers $a,b$ with $a$ odd, $b$ even,
and $a^2+b^2$ an integer square. The Gaussian integers $a+ib,a-ib$
are coprime: a common odd Gaussian prime would divide both $a,b$,
while $1+i$ does not divide their odd norm. Unique factorization in
$\Z[i]$ therefore makes $a+ib$ a unit times a square. Here unique
factorization follows from the norm-Euclidean algorithm: rounding both
coordinates of a complex quotient gives a remainder whose norm is
at most half that of the divisor. The unit is
not $\pm i$, because a Gaussian square of odd norm has odd real
part and even imaginary part. The remaining units $\pm1$ are
Gaussian squares. Writing $d_j=m_jt_j^2$ and absorbing $t_j$ gives
$U_j=m_jq_j^2$. Taking the positive square root of its norm proves
$A_j=m_j\Nm(q_j)$. Finally $U_j/L$ has squareclass $m_j/L$, which
equals $Lm_j$. For the converse, taking norms in the four displayed
root equations gives the four distances
$m_j\Nm(q_j)/L\in\Q_{>0}$.
\end{proof}

The product
\begin{equation}\label{eq:int:kernel}
 n=m_0m_1m_2m_3
\end{equation}
is thus a positive odd squarefree integer, prime to $L$, supported on
primes $1$ and $7$ modulo $8$. Equation
\eqref{eq:int:gaussian-classes} records both its corner allocation
and the single common factor $L$.

\begin{proposition}[The common squarefree kernel]\label{prop:int:kernel-product}
Put $t=r_0r_3+r_1r_2$. Then $t>1$ and
\[
 [t^2-1]=[n]\quad\text{in }\Q^*/\Q^{*2}.
\]
Consequently the Gaussian corner kernel \eqref{eq:int:kernel} equals
$\sf(t^2-1)$.
\end{proposition}
\begin{proof}
Write $z=x+iy$, and set
\[
 A=z(z-1-i),\qquad B=(z-1)(z-i),\qquad
 a=r_0r_3,\qquad b=r_1r_2.
\]
Then $B-A=i$, $A\overline A=a^2$ and $B\overline B=b^2$.
The triangle inequality gives $a+b\geq1$. Equality forces $A$ and $B$
to lie on the imaginary axis, and hence
\[
 \operatorname{Re}(A)=(x-y)(x+y-1)=0.
\]
Lemma~\ref{lem:int:symmetry} excludes equality, so $t>1$.
Expanding with the two norm identities gives
\[
 (bA+aB)^2
 =AB\bigl(A\overline B+B\overline A+2ab\bigr)
 =AB\bigl((a+b)^2-|A-B|^2\bigr)
 =(t^2-1)AB.
\]
Multiplication of the four Gaussian identities
\eqref{eq:int:gaussian-roots} gives
$AB=n(q_0q_1q_2q_3/L^2)^2$. Thus $(t^2-1)/n$ is a square
in $\Q(i)$. A positive rational number which is a square in
$\Q(i)$ is a square in $\Q$: if $(u+iv)^2$ is positive rational,
then $uv=0$ and positivity forces $v=0$. This proves the claim.
\end{proof}

The common factor also occurs in the unit restrictions at ramified
split primes. For $p\equiv1\pmod8$, choose either square root
$\iota$ of $-1$ in $\F_p$ and define
\[
 \epsilon(p)=\leg{1+\iota}{p}.
\]
This is independent of the choice: replacing $\iota$ by $-\iota$
multiplies $1+\iota$ by $-\iota$, a square in $\F_p^*$.
Equivalently,
\begin{equation}\label{eq:int:quartic-sign}
 \epsilon(p)=2^{(p-1)/4}(-1)^{(p-1)/8}\in\{1,-1\},
\end{equation}
where the power is evaluated in $\F_p$ and the result identified
with its sign.

\begin{proposition}\label{prop:int:ramified-units}
If $p\equiv1\pmod8$ divides $m_j$, then
\[
 \leg{Lm_k}{p}=1
\]
for either adjacent corner $k$, and
\[
 \leg{Lm_{j^{\mathrm{op}}}}p=\epsilon(p)
\]
for the opposite corner $j^{\mathrm{op}}$.
\end{proposition}

\begin{proof}
Use the reduction homomorphism $\Z[i]\longrightarrow\F_p$ given by
$i\mapsto\iota$. From \eqref{eq:int:gaussian-roots},
$X+iY\equiv Lc_j\pmod p$, and therefore for $k\ne j$
\[
 m_kq_k^2\equiv L(c_j-c_k)\pmod p.
\]
All factors on the right are nonzero. Adjacent differences are
$\pm1$ or $\pm i$, all squares modulo $p$. Opposite differences are
these units times $1+i$. Taking quadratic characters proves the
assertions. For \eqref{eq:int:quartic-sign}, use
$(1+\iota)^2=2\iota$ and raise to $(p-1)/4$.
\end{proof}

\subsection{Signed radii and denominator growth}

The integral equations also connect divisibility to the real size of
the square. Set
\[
 d=2+\sqrt2.
\]
For every strict point,
\begin{equation}\label{eq:int:total-bound}
 A_0+A_1+A_2+A_3<dL.
\end{equation}
Indeed, the sum of the distances to the four vertices is a convex
function on the square and equals $d$ at each vertex. A strict point
is a convex combination of all four vertices with positive weights.
The triangle inequality for the distance to $P_0$ is strict in this
combination, because the two adjacent nonzero vertex vectors are not
parallel. This proves \eqref{eq:int:total-bound}.

Choose arbitrary signs and write $a,b,c,e$ for the signed versions
of $A_0,A_1,A_2,A_3$. Put
\[
 S=a+b+c+e,\quad T=a-b-c+e,\quad
 P=ae-bc,\quad Q=ae+bc.
\]
The four-distance equations give the identities
\begin{equation}\label{eq:int:signed-identities}
 a^2+e^2=b^2+c^2,\qquad
 PQ=-L^2(L-2X)(L-2Y),\qquad ST=2P.
\end{equation}
The second follows by expanding $A_0^2A_3^2-A_1^2A_2^2$;
the last follows from the first. Lemma~\ref{lem:int:symmetry} shows
that the product on the right is nonzero. Thus $S,T,P,Q$ are all
nonzero, for every sign choice. Moreover
\begin{equation}\label{eq:int:signed-bound}
 |T|\le A_0+A_1+A_2+A_3<dL.
\end{equation}

\begin{theorem}[Prime-power comparisons]\label{thm:int:prime-comparisons}
Let $K=2^{v_2(L)}$. Then
\begin{equation}\label{eq:int:two-bound}
 K^2<\frac{2+\sqrt2}{4}\,L.
\end{equation}
If $p$ is an odd prime with $p\mid L$ and $p\mid XY$, and
$P_p=p^{v_p(L)}$, then
\begin{equation}\label{eq:int:odd-bound}
 P_p^2<\frac{2+\sqrt2}{2p}\,L.
\end{equation}
In particular
\begin{equation}\label{eq:int:three-bound}
 3^{2v_3(L)}<\frac{2+\sqrt2}{6}\,L.
\end{equation}
\end{theorem}

\begin{proof}
For the dyadic assertion use $X$ odd and $4\mid Y$. Write
$h=v_2(L)\ge3$ and $j=v_2(L-2Y)\ge3$; the latter is finite by
Lemma~\ref{lem:int:symmetry}. Choose the signs of $a,b,c,e$ so that
all four are $1$ modulo $4$. Every pair sum then has depth one.
The horizontal and vertical squared differences give
\[
 v_2(a-b)=v_2(c-e)=h,\qquad
 v_2(a-c)=v_2(b-e)=h+j-1.
\]
Consequently $a,b,c,e$ are congruent modulo $2^h$. It follows that
\[
 S\equiv4a\pmod{2^h},\qquad
 Q\equiv2a^2\pmod{2^h},
\]
so $v_2(S)=2$ and $v_2(Q)=1$. Since $v_2(L-2X)=1$, the identities
\eqref{eq:int:signed-identities} yield
\[
 v_2(P)=2h+j,\qquad v_2(T)=2h+j-1\ge2h+2.
\]
The nonzero integer $T$ therefore has $|T|\ge4K^2$.
Equation~\eqref{eq:int:signed-bound} proves \eqref{eq:int:two-bound}.

For the odd-prime assertion, exchange the coordinates independently
if needed so that $p\nmid X$ and $p\mid Y$. This is possible by
primitivity and the hypothesis $p\mid XY$. The four $A_i$ are
$p$-adic units and their squares have the same residue. Choose their
signs so that $a,b,c,e$ have the same nonzero residue modulo $p$.
Then $S\equiv4a$ and $Q\equiv2a^2\pmod p$ are units.
Set $h=v_p(L)$ and $j=v_p(L-2Y)\ge1$, again finite. The same
identities now give the exact depth
\[
 v_p(T)=2h+j\ge2h+1.
\]
Every $A_i$ is odd, so the signed sum $T$ is even. Hence
\[
 2pP_p^2\le2p^{2h+j}\le |T|<dL,
\]
which proves \eqref{eq:int:odd-bound}. At $p=3$, the required
divisibility $3\mid XY$ was proved in
Proposition~\ref{int:primitive}, so \eqref{eq:int:three-bound}
follows.
\end{proof}

Writing $L=Km$ with $m$ odd, \eqref{eq:int:two-bound} takes the
linear form
\[
 K<\frac{2+\sqrt2}{4}\,m.
\]
For odd $p$ as in the theorem, writing $L=P_ps$ similarly gives
$P_p<(2+\sqrt2)s/(2p)$. These comparisons apply to the complete
prime powers in the denominator.

\begin{corollary}\label{cor:int:six-smooth}
The least coordinate denominator of a strict rational-distance point
has a prime divisor at least $5$.
\end{corollary}

\begin{proof}
If $L$ were supported on $2$ and $3$, write $L=KP$ with
$K=2^{v_2(L)}$ and $P=3^{v_3(L)}$. The preceding theorem gives
\[
 K^2<dL/4,\qquad P^2<dL/6.
\]
Multiplication and $KP=L$ yield $1<d^2/24$. But
$d^2=6+4\sqrt2<24$, a contradiction.
\end{proof}

The same signed identities give a precise statement in the other
$5$-adic direction, where all four radii have nontrivial content.

\begin{proposition}[The unit-coordinate direction at $5$]
\label{prop:int:five-depth}
Suppose $5\mid L$ and $5\nmid XY$. Write $L=5^h m$ with $5\nmid m$.
There is an integer $s\ge1$ such that
\[
 v_5(A_i)=s\quad(0\le i\le3),\qquad h>2s.
\]
Some signed sum $T$ of the four $A_i$ has exact depth $2h-3s$,
and
\begin{equation}\label{eq:int:five-bound}
 5^{h-3s}<\frac{2+\sqrt2}{2}\,m.
\end{equation}
\end{proposition}

\begin{proof}
Nonzero squares modulo $5$ are $1$ and $4$. Since $X,Y$ are units
and $X^2+Y^2$ is a square, this sum must be zero modulo $5$.
Thus $Y/X\equiv\pm2\pmod5$, and all four $A_i$ are divisible by $5$.
Let $s=\min_i v_5(A_i)$. The adjacent squared differences have
exact depth $h$, so $h\ge2s$.

If $h=2s$, division of the four squared-distance equations by
$5^{2s}$ gives residues
\[
 c,\quad c+u,\quad c+v,\quad c+u+v,
 \qquad u,v\ne0,\quad v/u=\pm2
\]
in $\F_5$. These four residues are distinct, whereas $\F_5$ has
only three square residues. Therefore $h>2s$. Every squared
difference has depth at least $h$, and hence all four radii have
depth exactly $s$.

Divide the four radii by $5^s$ and choose signs so that the resulting
units $a,b,c,e$ have the same residue modulo $5$. Their sums
$S=a+b+c+e$ and $Q=ae+bc$ are units. The determinant identity becomes
\[
 (ae-bc)(ae+bc)=-\frac{L^2(L-2X)(L-2Y)}{5^{4s}}.
\]
Both $L-2X,L-2Y$ are units, so $ae-bc$ and
$a-b-c+e$ have exact depth $2h-4s$. Multiplying the latter by
$5^s$ gives an actual signed sum of the integer radii with exact
depth $2h-3s$. It is nonzero and even. Its absolute value is at
least $2\cdot5^{2h-3s}$ and less than $dL$, which proves
\eqref{eq:int:five-bound}.
\end{proof}
